\section*{第 \ref{ch:b1:cross-sections} 章——剖面}

\begin{solution}{exo:b1:cross-sections:1}
在 1:25\,000 的比例尺下，地图上的 \qty{1}{cm} 对应地面 \qty{250}{m}。
竖直比例尺为 \qty{1}{cm} 对应 \qty{100}{m}，因此竖直比例尺是水平比例尺的
$250/100 = 2.5$ 倍：$e = 2.5$。
\end{solution}

\begin{solution}{exo:b1:cross-sections:2}
$\gamma = \ang{90} - \ang{60} = \ang{30}$，且 $\tan\beta = \tan\ang{40}
\cos\ang{30} = 0.8391\times0.8660 = 0.7267$：$\beta = \ang{36.0}$，岩层
向剖面东端下倾。在\omterm{def:b1:spherical-astronomy:horizontal}{方位角}
\ang{000} 的剖面中，$\gamma = \ang{90}$ 且 $\beta = 0$：该剖面沿
\omterm{def:b1:geological-maps:attitude}{走向}延伸，岩层呈水平。
\end{solution}

\begin{solution}{exo:b1:cross-sections:3}
$t = h\cos\delta = 120\times\cos\ang{12.16} = 120\times0.9776 =
\qty{117.3}{m}$。
\end{solution}

\begin{solution}{exo:b1:cross-sections:4}
$\tan\ang{12} = 0.2126$。取 $e = 2$：$\arctan 0.4251 = \ang{23.0}$；取
$e = 3$：$\arctan 0.6377 = \ang{32.5}$；取 $e = 5$：$\arctan 1.0628 =
\ang{46.7}$。\omterm{def:b1:geological-maps:attitude}{倾角}的增长远小于 $e$ 本身，且永远达不到
\ang{90}。
\end{solution}

\begin{solution}{exo:b1:cross-sections:5}
(a) $t = 300\sin\ang{25} = \qty{126.8}{m}$。(b) $t = 300\sin\ang{35} =
\qty{172.1}{m}$。(c) $t = 300\sin\ang{15} = \qty{77.6}{m}$。同样宽度的\omterm{def:b1:geological-maps:geological-map}{露头}
可对应相差两倍以上的厚度：
必须始终考虑地面的坡度。
\end{solution}

\begin{solution}{exo:b1:cross-sections:6}
设\omterm{def:b1:geological-maps:attitude}{倾向}为 $\alpha$、\omterm{def:b1:geological-maps:attitude}{倾角}为 $\delta$，由定理可得，在
南北向剖面（$\theta = 0$）中 $\tan\ang{20} = \tan\delta\cos\alpha$，
在东西向剖面（$\theta = \ang{90}$）中 $\tan\ang{30} =
\tan\delta\cos(\alpha - \ang{90}) = \tan\delta\sin\alpha$。两式相除，
$\tan\alpha = 0.5774/0.3640 = 1.586$，$\alpha = \ang{057.8}$；两式
平方相加，$\tan\delta = \sqrt{0.3640^2 + 0.5774^2} = 0.6825$，$\delta =
\ang{34.3}$。两个\omterm{def:b1:cross-sections:apparent-dip}{视倾角}都小于真\omterm{def:b1:geological-maps:attitude}{倾角}，
这是必然的。
\end{solution}

\begin{solution}{exo:b1:cross-sections:7}
$100/\tan\ang{12.16} = 100/0.2155 = \qty{464.1}{m}$。在地图上，
东翼的\omterm{def:b1:cross-sections:structure-contour}{构造等高线}确实相距略小于半
千米，且等距、平行：该翼为一平面。
\end{solution}

\begin{solution}{exo:b1:cross-sections:8}
B1：$524.02 - 488.98 = \qty{35.04}{m}$；B3：$529.98 - 494.93 =
\qty{35.05}{m}$。在 B2 之上，页岩底面即石灰岩顶面，
石灰岩顶面高出其底面 \qty{180}{m}：$280.16 + 180 = \qty{460.16}{m}$，
高于地面（\qty{380.64}{m}），因此该处已被侵蚀。
在 B1 与 B2 之间，界线上升 $460.16 - 35.04 = \qty{425.12}{m}$，
水平距离为 \qty{2000}{m}：坡度为 $0.2126$，\omterm{def:b1:cross-sections:apparent-dip}{视倾角}为 \ang{12.0}。B1
与 B3 位于 B2 两侧相同距离处，它们测得的各界线的
高程均在 \qty{1}{cm} 之内相同：该\omterm{def:b1:geological-maps:fold}{褶皱}关于竖直
轴面 $x = \qty{3.5}{km}$ 对称，两翼\omterm{def:b1:geological-maps:attitude}{倾角}相等，且两者之间没有
断层切割剖面。
\end{solution}

\begin{solution}{exo:b1:cross-sections:9}
\omterm{def:b1:cross-sections:throw}{水平断距} $= 150/\tan\ang{50} = 150/1.1918 = \qty{125.9}{m}$。在
\omterm{def:b1:spherical-astronomy:horizontal}{方位角}为 \ang{150} 的剖面中，$\gamma = \ang{30}$，且 $\tan\beta = 1.1918\times0.8660 =
1.0321$：断层呈现\omterm{def:b1:cross-sections:apparent-dip}{视倾角} \ang{45.9}。水平岩层的截止线
在该迹线上竖直相距 \qty{150}{m}，因此
沿剖面相距 $150/1.0321 = \qty{145.3}{m}$，大于\omterm{def:b1:cross-sections:throw}{水平断距}，
因为剖面斜交于断层。
\end{solution}

\begin{solution}{exo:b1:cross-sections:10}
$\cos\gamma = \tan\ang{20}/\tan\ang{40} = 0.3640/0.8391 = 0.4338$，故
$\gamma = \ang{64.3}$：\omterm{def:b1:spherical-astronomy:horizontal}{方位角}为 $\alpha \pm \ang{64.3}$ 的剖面
（以及相反方向上岩层上升的剖面）。\omterm{def:b1:cross-sections:apparent-dip}{视倾角}为
\ang{40.0}、\ang{39.6}、\ang{38.3} 和 \ang{36.0}，分别对应 $\gamma = 0$、\ang{10}，
\ang{20}、\ang{30}；而 \ang{22.8}、\ang{16.0}、\ang{8.3} 和 $0$ 则对应
$\gamma = \ang{60}$、\ang{70}、\ang{80}、\ang{90}。事实上，对
$\tan\beta = \tan\delta\cos\gamma$，求导可得
$\dfrac{\mathrm d\beta}{\mathrm d\gamma} =
-\dfrac{\tan\delta\sin\gamma}{1 + \tan^2\delta\cos^2\gamma}$，其
在 $\gamma = 0$ 时为零，取值 $-\tan\delta$ 的是 $\gamma = \ang{90}$ 处：稍偏离\omterm{def:b1:geological-maps:attitude}{倾向}方向的
剖面几乎显示真\omterm{def:b1:geological-maps:attitude}{倾角}，而接近\omterm{def:b1:geological-maps:attitude}{走向}时，
\omterm{def:b1:spherical-astronomy:horizontal}{方位角}每变化一度，\omterm{def:b1:cross-sections:apparent-dip}{视倾角}几乎随之变化一度。
\end{solution}

\begin{solution}{exo:b1:cross-sections:11}
岩层向东倾，因此沿层面法线向下，其
相对竖直方向向西倾斜 \ang{12.16}，即朝向 \ang{260.7}。
(a) 竖直钻孔与法线的夹角为 $\psi = \ang{12.16}$：
$176.0/\cos\ang{12.16} = \qty{180.0}{m}$，即\omterm{def:b1:cross-sections:thickness}{铅直厚度}。(b) 一个
朝向 \ang{260.7} 倾伏 \ang{60} 的钻孔，与竖直方向的夹角为 \ang{30}，
且与法线位于同一侧：$\psi = \ang{30} - \ang{12.16} =
\ang{17.84}$，$\cos\psi = 0.9519$，岩芯长为 $176.0/0.9519 =
\qty{184.9}{m}$。(c) 朝向 \ang{080.7} 时则倾斜于另一侧：
$\psi = \ang{30} + \ang{12.16} = \ang{42.16}$，$\cos\psi = 0.7413$，岩芯长为
$176.0/0.7413 = \qty{237.4}{m}$。岩芯最短，等于
$t$，此时钻孔沿法线，朝向
\ang{260.7} 倾伏 \ang{77.84}。
\end{solution}

\begin{solution}{exo:b1:cross-sections:12}
在东西向竖直面内，角度自东向量取，
上部岩层向东上倾，\omterm{def:b1:geological-maps:attitude}{倾角}为 \ang{20}；下部岩层向东下倾，
\omterm{def:b1:geological-maps:attitude}{倾角}为 \ang{70}。绕\omterm{def:b1:geological-maps:attitude}{走向}将\omterm{def:b1:geological-maps:geological-map}{露头}
旋转 \ang{20}（东侧向下），使上部岩层水平，并使
下部岩层转为 $\ang{70} + \ang{20} = \ang{90}$，即在
沉积上部岩层时，下部岩层是直立的。按时间顺序，从最老的事件开始：下部岩层
平放沉积；随后被翻转至直立，并抬升；
其边缘被侵蚀成一个面；上部岩层沉积覆盖于该
面之上；整体向西倾斜 \ang{20}；侵蚀
切割出今天的地表。这正是赫顿在锡卡尔
角\omterm{def:b1:geological-maps:geological-map}{露头}所记录的历史。
\end{solution}

\begin{solution}{pb:b1:cross-sections:1}
\textbf{1.} 隧道上方的岩层厚度为 $z - 200$，在
$x = \qty{1}{km}$ 处最大，$556 - 200 = \qty{356}{m}$；在 $x =
\qty{9}{km}$ 处最小，$339 - 200 = \qty{139}{m}$。

\textbf{2.} 图上 \qty{1}{km} 的高程量得 \qty{4.2}{cm}，\qty{1}{km}
的水平距离量得 \qty{1.4}{cm}：$e = 4.2/1.4 = 3$。

\textbf{3.} $380.64 - 220.48 = \qty{160.16}{m}$。

\textbf{4.} 在 B2 处，石灰岩底面高程为 $380.64 - 100.48 =
\qty{280.16}{m}$，砂岩底面为 \qty{160.16}{m}：隧道
在 \qty{200}{m} 高程处位于砂岩中。

\textbf{5.} B1：$524.02 - 668.98 = \qty{-144.96}{m}$；B3：$529.98 - 674.93 =
\qty{-144.95}{m}$。

\textbf{6.} 由西向东依次为白垩、页岩、石灰岩、砂岩（位于\omterm{def:b1:geological-maps:fold}{背斜}核部）、石灰岩、页岩，
以及白垩，直至
$x = \qty{9}{km}$ 处，途经\omterm{def:b1:geological-maps:fold}{向斜}并穿越断层。

\textbf{7.} $\gamma = \ang{80.7} - \ang{90} = \ang{-9.3}$；$\tan\beta =
\tan\ang{12.16}\cos\ang{9.3} = 0.2155\times0.9869 = 0.2126$：$\beta =
\ang{12.0}$，向东下倾。

\textbf{8.} $(280.16 + 144.95)/2000 = 425.11/2000 = 0.2126$：\ang{12.0}。

\textbf{9.} $\arctan(3\times0.2126) = \arctan 0.6377 = \ang{32.5}$。

\textbf{10.} $\gamma = \ang{120} - \ang{90} = \ang{30}$；$\tan\beta =
\tan\ang{60}\cos\ang{30} = 1.732\times0.866 = 1.500$：$\beta = \ang{56.3}$。

\textbf{11.} 两翼\omterm{def:b1:geological-maps:attitude}{倾向}为 \ang{080.7} 或 \ang{279.3}，
与剖面线仅相差 \ang{9.3}，且 $\cos\ang{9.3} = 0.987$；断层\omterm{def:b1:geological-maps:attitude}{倾向}为 \ang{120}，
与剖面线相差 \ang{30}，且 $\cos\ang{30} = 0.866$。

\textbf{12.} 断层面为 $z = 250 - \tan\ang{60}\,[(x - 4300)
\sin\ang{120} + (y - 1000)\cos\ang{120}]$。取 $y = 6000$，则 $(y - 1000)
\cos\ang{120} = -2500$，故 $(x - 4300)\sin\ang{120} = 2500 - (z - 250)/
\tan\ang{60}$，且 $1/\tan\ang{60} = \tan\ang{30}$。

\textbf{13.} 当 $z = 200$ 时：$x = 4300 + (2500 + 50\times0.5774)/0.8660 =
4300 + 2528.9/0.8660 = \qty{7220}{m}$，即距 C 点 \qty{7.22}{km}。

\textbf{14.} \omterm{def:b1:cross-sections:throw}{水平断距} $= 200/\tan\ang{60} = \qty{115.5}{m}$；沿
C--D 方向的间距 $= 200/1.500 = \qty{133.3}{m}$。

\textbf{15.} 断层向东南东倾斜，为正断层：位于东侧的上盘下降。
在隧道高程处，断层两侧均为白垩：断层西侧为
\omterm{def:b1:geological-maps:fold}{向斜}中的白垩；东侧为
被下降 \qty{200}{m} 的白垩。

\textbf{16.} 隧道穿越断层时始终位于白垩中，因此
掌子面处岩性不变，无法据此察觉断层；只有破碎带，
或许还有地下水，才能显示断层。工程师可依据剖面预测
其位置（此处为 \qty{7.22}{km}），在掘进机接近时提前
钻探超前探孔，并为破碎带准备支护。

\textbf{17.} $668.98 - 488.98 = \qty{180.00}{m}$。

\textbf{18.} $t = 180\cos\ang{12.16} = \qty{176.0}{m}$。

\textbf{19.} 在\omterm{def:b1:geological-maps:fold}{背斜}轴顶（$x = 3500$）处，砂岩底面高程为
\qty{160.16}{m}（见第 3 题）；两翼均向下，
\omterm{def:b1:cross-sections:apparent-dip}{视倾角}为第 7 题中的 \ang{12.0}，即
距轴线每米水平距离下降 $\tan\ang{12.0}$。

\textbf{20.} 石灰岩底面高于砂岩底面 \qty{120}{m}，位于
$280.16 - u\tan\ang{12.0}$ 处，其中 $u = |x - 3500|$，其顶面
再高 \qty{180}{m}，位于 $460.16 - u\tan\ang{12.0}$。取 $\tan\ang{12.0} =
0.21256$：当 $u > 80.16/0.21256 = 377.1$ 时底面低于 \qty{200}{m}，
当 $u < 260.16/0.21256 = 1224.0$ 时顶面高于 \qty{200}{m}。

\textbf{21.} $x = 3500 \pm 377.1$ 与 $3500 \pm 1224.0$：隧道位于
石灰岩中的范围为 \qty{2276}{m} 至 \qty{3123}{m}，以及 \qty{3877}{m} 至
\qty{4724}{m}。

\textbf{22.} 沿水平线方向，石灰岩的底面与顶面竖直相距 \qty{180}{m}，
二者均以 $\tan\beta$ 倾斜，沿水平线穿过的水平间距为 $h/\tan\beta$：
每翼为 $180/0.21256 = \qty{846.8}{m}$，
与舍入后的 $1224.0 - 377.1 = 846.9$ 一致。

\textbf{23.} $2\times846.8 = \qty{1693.7}{m}$：隧道在石灰岩中掘进
\textbf{约 \qty[text-series-to-math=true]{1694}{m}}，超过其 \qty{8}{km} 全长的五分之一，
分为两段，每段 \qty{847}{m}，位于
\omterm{def:b1:geological-maps:fold}{背斜}砂岩核部两侧。
\end{solution}
